% Détermination pH-métrique d'une constante d'équilibre
Une solution d'acide méthanoïque de concentration molaire $C=\qty{1,0e-2}{\M}$ a
un pH égal à 2,9 à $\qty{25}{\degreeCelsius}$.
\begin{enumerate}
    \item Donner l'équation de la réaction de l'acide méthanoïque avec l'eau.
    \item Montrer que la réaction de l'acide méthanoïque avec l'eau n'est pas
          totale.
    \item Calculer le quotient de réaction associé à l'équation de cette
          réaction dans l'état d'équilibre du système.
    \item En déduire la valeur de la constante d'équilibre associée à cette
          équation.
\end{enumerate}

